\documentclass[11pt,a4paper]{article}
\usepackage[T1]{fontenc}
\usepackage[utf8]{inputenc}
\usepackage{mathptmx}
\AtBeginDocument{\let\hbar\relax
  \DeclareMathSymbol{\hbar}{\mathord}{AMSb}{"7E}}
\usepackage{amsmath,amssymb,amsthm}
\usepackage{amsfonts}
\usepackage{geometry}
\usepackage{microtype}
\usepackage{hyperref}
\usepackage{booktabs}
\usepackage{enumitem}
\usepackage{natbib}
\usepackage{titlesec}

\geometry{top=1.1in,bottom=1.1in,left=1.15in,right=1.15in}

\titleformat{\section}{\normalfont\bfseries\large}{\thesection}{1em}{}
\titleformat{\subsection}{\normalfont\bfseries\normalsize}{\thesubsection}{1em}{}
\titleformat{\subsubsection}{\normalfont\bfseries\itshape\normalsize}{\thesubsubsection}{1em}{}
\titlespacing*{\section}{0pt}{14pt}{4pt}
\titlespacing*{\subsection}{0pt}{10pt}{3pt}
\titlespacing*{\subsubsection}{0pt}{8pt}{2pt}

\theoremstyle{plain}
\newtheorem{theorem}{Theorem}[section]
\newtheorem{proposition}[theorem]{Proposition}
\newtheorem{corollary}[theorem]{Corollary}
\theoremstyle{definition}
\newtheorem{definition}[theorem]{Definition}
\newtheorem{conjecture}[theorem]{Conjecture}
\theoremstyle{remark}
\newtheorem{remark}[theorem]{Remark}

\hypersetup{colorlinks=true,linkcolor=black,citecolor=black,urlcolor=black,
  pdftitle={Relational Actualism: Dark Matter as Actualization-Density Topology (RADM)},
  pdfauthor={Joshua F. Sandeman}}

\newcommand{\Pact}{\mathcal{P}_{\mathrm{act}}}
\newcommand{\Tact}{T^{\mu\nu}_{\mathrm{RA}}}
\newcommand{\ARA}{A_{\mathrm{RA}}}
\newcommand{\GG}{\mathcal{G}}
\newcommand{\AP}{\mathrm{AP}}
\newcommand{\MB}{\mathrm{MB}}
\newcommand{\ket}[1]{|#1\rangle}

% ─────────────────────────────────────────────────────────────────────────────
\begin{document}

% ── Title block ─────────────────────────────────────────────────────────────
\begin{center}
  {\LARGE\bfseries Relational Actualism: Dark Matter as\\
  Actualization-Density Topology (RADM)\\[4pt]
  \large\itshape Flat Rotation Curves, the Bullet Cluster,\\
  and the Topological Boundary Mechanism\par}
  \vspace{14pt}
  {\large Joshua F.\ Sandeman$^{1*}$\par}
  \vspace{4pt}
  {\normalsize $^{1}$Independent Researcher, Salem, OR, USA.\par}
  {\normalsize $^{*}$sansuikyo@gmail.com\par}
  \vspace{8pt}
  {\small\textit{Preprint available at}~\href{https://doi.org/10.5281/zenodo.19198513}{doi.org/10.5281/zenodo.19198513}\par}
  \vspace{14pt}
\end{center}
\noindent\textbf{Ethics statement.}
This article does not contain any studies with human participants or animals.
\medskip
\noindent\textbf{AI use disclosure.}
Claude (Anthropic, \texttt{claude-sonnet-4-6}) and Gemini (Google DeepMind)
were used in manuscript preparation. The author is solely responsible for all
scientific content and conclusions.
\medskip
\begin{abstract}
Relational Actualism (RA) prohibits causally silent dark matter candidates
structurally: a 100~GeV WIMP generates $10^{25}$--$10^{38}$ times less
actualization flux than equivalent baryonic matter, falling categorically below
the threshold for gravitational lensing. This paper demonstrates that the
observed phenomenology attributed to dark matter is accounted for by the
actualization-density field of RA, without invoking any new particle species,
across four distinct observational regimes.

\emph{Rotation curves.} The 2D cylindrical Laplacian of $\ln\lambda$
gives $\Sigma_\lambda \propto 1/R$ for exponential disk profiles, producing
flat rotation curves with $v_{\mathrm{flat}}^2 = \xi/(2r_d)$ (RA-native
derivation). In the sparse halo regime, the Durhuus-Jonsson-Wheater theorem
for Lorentzian graphs gives $\Phi \propto \ln r$, from which the
Tully-Fisher relation $v_{\mathrm{flat}}^4 \propto M$ follows by causal
flux conservation.

\emph{Bullet Cluster.} The $A_{\mathrm{RA}}$(stellar)/$A_{\mathrm{RA}}$(gas)
ratio of $\sim 10^8$ causes lensing to trace the stellar component
dynamically, without requiring a dark matter halo or leaving a trailing
gravitational wake.

\emph{Galaxy-galaxy lensing.} The DJW logarithmic potential gives the
convergence profile $\kappa(b) = v_{\mathrm{flat}}^2/(4G\Sigma_{\mathrm{cr}}b)$,
cutting off at the baryonic edge $R_{\mathrm{baryon}} \sim 150$--$300$~kpc.
The Einstein radius is $\theta_E \approx 0.47$~arcsec for a Milky-Way
analog, consistent with observed strong-lensing systems.

\emph{Hubble tension.} A bandwidth-partition model gives
$H_{\mathrm{local}} \approx 73.1$~km\,s$^{-1}$\,Mpc$^{-1}$ from the KBC
void. New falsifiable prediction: the Eridanus supervoid
($\delta \approx -0.30$) should give $H \approx 75.9$~km\,s$^{-1}$\,Mpc$^{-1}$.

Five hard walls are precisely scoped for collaborators.

\end{abstract}

% ─────────────────────────────────────────────────────────────────────────────
\section{Introduction: The Dark Matter Problem in Relational Actualism}
\label{sec:intro}
% ─────────────────────────────────────────────────────────────────────────────

The galactic rotation curve problem, the Bullet Cluster lensing offset, and the
large-scale structure formation requirements of the standard $\Lambda$CDM model
are three of the strongest constraints in observational cosmology. The standard
account invokes a non-baryonic cold dark matter (CDM) component---a particle
species that gravitates but does not interact electromagnetically---to satisfy all
three simultaneously.

The companion paper RAGC~\citep{Sandeman2026b} establishes that this account is
structurally incompatible with Relational Actualism. In RA, the metric is sourced
by actualization flux $\lambda(x)$, not rest-mass density. A WIMP undergoes at
most one on-shell interaction every $10^{18}$ universe lifetimes; it contributes
negligible actualization flux and therefore negligible metric sourcing. This is the
WIMP Prohibition (RAGC Prediction~4). The prohibition extends to axions, sterile
neutrinos, and any species with negligible electromagnetic and strong cross-sections.

This paper asks the complementary question: if causally silent particles cannot
generate the dark matter signatures, what does? We propose that the answer lies in
the topological structure of the actualization-density field itself---specifically,
in two mechanisms that follow from the same RA framework:

\begin{enumerate}[noitemsep]
  \item \textbf{The $\nabla^2(\ln\lambda)$ rotation curve mechanism:} a gradient
  correction to the metric sourcing equation that produces flat rotation curves
  for exponential disk profiles at natural coupling order.
  \item \textbf{The accumulated causal depth mechanism:} the distinction between
  instantaneous flux $\lambda(x,t)$ and accumulated causal depth $A_{\mathrm{RA}}$
  as the metric source, which resolves the Bullet Cluster lensing offset without
  new parameters.
\end{enumerate}

Both mechanisms are conjectural. Neither constitutes a complete replacement for
CDM until the hard walls identified in Section~\ref{sec:hardwalls} are resolved.
A key conceptual shift informs this paper: RADM does not compete with GR
but predicts exactly where GR fails. In dense baryonic regions ($\mu \gg 1$),
the causal graph IS 3-dimensional and GR is recovered as the unique
consistent macroscopic description (Lovelock's theorem~\citep{Sandeman2026b}).
In sparse galactic halos ($\mu < 1$), the graph IS locally 2-dimensional
and GR's assumption of a uniform 4D manifold breaks down. Dark matter is not
needed --- the geometry itself changes. This paper presents the derivations
as far as they currently reach, names the hard walls precisely, and identifies
the collaborator targets that would resolve them.

% ─────────────────────────────────────────────────────────────────────────────
\section{The Unified Metric Sourcing Equation}
\label{sec:sourcing}
% ─────────────────────────────────────────────────────────────────────────────

The analysis of the WIMP Prohibition in RAGC reveals a deeper structural result.
The weak-field metric sourcing equation in RA separates into two terms that
standard GR conflates because for ordinary matter in equilibrium they are
proportional:

\begin{equation}
  \nabla^2\Phi = 4\pi G\bigl[\rho_A + \rho_\lambda\bigr],
  \label{eq:unified}
\end{equation}
where:
\begin{align}
  \rho_A(x) &= \frac{\hbar\langle\Delta E\rangle}{c^2}
    \cdot \frac{\ARA(x)}{V}, \label{eq:rhoA} \\
  \rho_\lambda(x) &= -\frac{\xi}{4\pi G}\,\nabla^2(\ln\lambda(x)).
    \label{eq:rhoL}
\end{align}

The first term $\rho_A$ is sourced by accumulated causal depth---the integral of
$\lambda$ over the past light cone $J^-(x)$. For ordinary matter in thermodynamic
equilibrium over cosmological time, $\ARA/V \approx \lambda \cdot t_{\mathrm{cosm}}
\cdot \langle\Delta E\rangle/c^2 \propto \rho_{\mathrm{matter}}$: this term
recovers standard GR, and the weak equivalence principle is respected for all
equilibrium species (see RAGC Derivation~6~\citep{Sandeman2026b} for the
WEP recovery argument).

The second term $\rho_\lambda$ is a dynamical correction proportional to the
Laplacian of the log of the actualization density. It vanishes when $\lambda$ is
spatially uniform---recovering standard GR in the equilibrium, homogeneous
limit---and becomes significant at boundaries between high-$\lambda$ and
low-$\lambda$ regions.

\begin{remark}[Physical motivation of the $\rho_\lambda$ term]
The quantity $\ln\lambda$ is proportional to the thermodynamic entropy of the
actualization field. Its Laplacian $\nabla^2(\ln\lambda)$ measures the local
divergence of actualization flow---positive where actualization is converging
(dense regions) and negative where it is spreading (sparse regions). The
$\rho_\lambda$ term encodes the topological tension at boundaries between
dense and sparse causal subgraphs.
\end{remark}

\paragraph{The four regimes where the terms separate.}
Standard GR is recovered everywhere both terms are proportional. They separate
observably only at:
\begin{enumerate}[noitemsep]
  \item \textbf{Galactic edges} ($\nabla\lambda$ large): $\rho_\lambda$ produces
  a boundary correction that flattens rotation curves (Section~\ref{sec:rotation}).
  \item \textbf{Intergalactic voids} ($\lambda \to 0$): $\rho_A$ and $\rho_\lambda$
  both suppressed, producing the differential expansion rate responsible for the
  Hubble tension (RAGC Prediction~1~\citep{Sandeman2026b}).
  \item \textbf{Transient high-$\lambda$ events} ($\lambda$ spikes but
  $\ARA$ has not accumulated): $\rho_A$ dominates over the gas, as in the
  Bullet Cluster (Section~\ref{sec:bullet}).
  \item \textbf{Causally silent species} ($\lambda \approx 0$ at all times):
  $\rho_A \approx 0$ and $\rho_\lambda \approx 0$; no metric sourcing
  (RAGC WIMP Prohibition~\citep{Sandeman2026b}).
\end{enumerate}

\paragraph{Covariant origin.}
Equation~\eqref{eq:unified} is no longer an ad-hoc phenomenological postulate.
As established in the companion paper RAGC~\citep{Sandeman2026b}, it emerges as
the strict static weak-field limit ($00$-component) of the fully covariant RA
field equation:
\begin{equation}
  G_{\mu\nu} = 8\pi G\!\left(T_{\mu\nu}^{(A)}
    + \Theta_{\mu\nu}^{(\lambda)}\right),
  \label{eq:covariant_field}
\end{equation}
where $\Theta_{\mu\nu}^{(\lambda)} = \xi[\nabla_\mu\nabla_\nu(\ln\lambda)
- g_{\mu\nu}\Box(\ln\lambda)]$ is the actualization tensor representing the
thermodynamic entropy flux of the graph-rewriting process. In the
non-relativistic regime where spatial gradients dominate time derivatives,
$G_{00} \approx \nabla^2\Phi$ and
$\Theta_{00}^{(\lambda)} \approx \xi\nabla^2(\ln\lambda)$, cleanly recovering
equation~\eqref{eq:unified}. Dimensional analysis gives
$\xi \sim r_d^2 v_{\mathrm{flat}}^2$ (natural order); the numerical value
$\xi/(G M_{\mathrm{disk}}) = 0.33$ for a Milky-Way-analog galaxy confirms
natural coupling. The astrophysical consequences explored in this paper are
therefore grounded in a first-principles covariant framework. The remaining
open problem --- deriving $\xi$ from the microscopic actualization entropy ---
is Hard Wall~1 (Section~\ref{sec:hardwalls}).

% ─────────────────────────────────────────────────────────────────────────────
\section{Flat Galactic Rotation Curves}
\label{sec:rotation}
% ─────────────────────────────────────────────────────────────────────────────

\subsection{The $\nabla^2(\ln\lambda)$ Mechanism}

A galaxy disk is a geometrically thin structure: the stellar actualization
density $\lambda(R)$ varies primarily in the radial direction, not in the
vertical direction on scales relevant to orbital dynamics.
The correct Laplacian to apply is therefore the \emph{2D cylindrical} operator:
\begin{equation}
  \nabla^2_\text{2D}(\ln\lambda)
  = \frac{1}{R}\frac{d}{dR}\!\left(R\,\frac{d\ln\lambda}{dR}\right).
  \label{eq:laplacian_2D}
\end{equation}
For an exponential disk profile $\lambda(R) = \lambda_0\,e^{-R/r_d}$:
\begin{equation}
  \frac{d\ln\lambda}{dR} = -\frac{1}{r_d},
  \qquad
  \frac{d}{dR}\!\left(R\cdot\left(-\frac{1}{r_d}\right)\right)
  = -\frac{1}{r_d},
  \qquad
  \nabla^2_\text{2D}(\ln\lambda) = -\frac{1}{R\,r_d}.
  \label{eq:lnlambda_laplacian}
\end{equation}
This yields an effective projected surface density:
\begin{equation}
  \Sigma_\lambda(R) = -\frac{\xi}{4\pi G}\,\nabla^2_\text{2D}(\ln\lambda)
  = \frac{\xi}{4\pi G\,R\,r_d}.
  \label{eq:Sigma_lambda_profile}
\end{equation}
The $1/R$ surface density profile integrates to a linearly growing enclosed mass:
\begin{equation}
  M_\lambda(<R) = \int_0^R 2\pi R'\,\Sigma_\lambda(R')\,dR'
  = \frac{\xi\,R}{2G\,r_d}.
  \label{eq:M_lambda}
\end{equation}
The circular velocity contribution is then:
\begin{equation}
  v^2_\lambda(R) = \frac{G\,M_\lambda(<R)}{R} = \frac{\xi}{2r_d} = \mathrm{const.}
  \label{eq:flat_curve}
\end{equation}
The $\nabla^2_\text{2D}(\ln\lambda)$ term in disk geometry produces a
\emph{perfectly flat} rotation curve, with
$v_{\mathrm{flat}} = \sqrt{\xi/(2r_d)}$.
The total rotation curve is:
\begin{equation}
  v_{\mathrm{total}}^2(R) = v_N^2(R) + \frac{\xi}{2r_d},
  \label{eq:v_total}
\end{equation}
where $v_N(R)$ is the standard Newtonian baryonic term.
The coupling constant calibrates to:
\begin{equation}
  \xi = 2\,v_{\mathrm{flat}}^2\,r_d.
  \label{eq:xi_calibration}
\end{equation}
For a Milky-Way analog ($M_{\mathrm{disk}} = 6\times10^{10}\,M_\odot$,
$r_d = 3.5$~kpc, $v_{\mathrm{flat}} = 220$~km~s$^{-1}$):
$\xi = 1.05\times10^{31}$~m$^3$~s$^{-2}$, and
$\xi/(GM_{\mathrm{disk}}) = 1.32$---natural order, no fine-tuning.
A numerical integration confirms flatness to within 4.3\% for $R > 5$~kpc.

\paragraph{Note on geometry.}
Applying the full 3D spherical Laplacian to the same exponential profile gives
$\nabla^2_\text{3D}(\ln\lambda) = -2/(r\,r_d)$, which produces a 3D volume
density $\rho_\lambda \propto 1/(r\,r_d)$. The enclosed mass from this profile
scales as $M \propto r^2$, giving a \emph{rising} rotation curve $v \propto r^{1/2}$.
The 3D spherical geometry is not appropriate for a thin disk; the 2D cylindrical
geometry correctly captures the disk structure and gives the flat result above.

\subsection{Comparison to MOND and Verlinde}

The $\nabla^2(\ln\lambda)$ mechanism is conceptually related to two
phenomenological approaches. Milgrom's MOND~\citep{Milgrom1983} introduces a
critical acceleration $a_0$ by hand to recover flat rotation curves below a
threshold. Verlinde's emergent gravity~\citep{Verlinde2011} models dark matter as
an elastic response of entanglement entropy in de Sitter space. The RA account
differs from both: MOND's $a_0$ is a free phenomenological parameter, while
the RA transition radius is set by the physical disk scale radius $r_d$. Verlinde
requires de Sitter space and entanglement entropy as primitives; RA derives the
effect from the causal graph. RA provides the fundamental physical mechanism; MOND
describes the phenomenological consequence.


% ─────────────────────────────────────────────────────────────────────────────
\subsection{Dimensional Reduction and the Tully-Fisher Limit}
\label{sec:dimensional_reduction}
% ─────────────────────────────────────────────────────────────────────────────

The $\nabla^2(\ln\lambda)$ mechanism of Section~\ref{sec:rotation}.1 is the
correct description in the dense baryonic regime where the causal DAG is
tightly woven and the continuum approximation holds. However, the Bianchi
constraint analysis reveals that the $\lambda$ profile saturates at large $r$
once the baryonic density vanishes, causing the rotation curve to turn over
rather than remain flat. This is not a failure of the actualization ontology;
it is the breakdown of the continuum approximation in the sparse galactic halo.

In discrete quantum gravity --- specifically Causal Dynamical
Triangulations and Causal Set Theory (see RAGC~\citep{Sandeman2026b} for references) ---
it is a rigorously established phenomenon that sparse causal graphs undergo
\emph{dynamical dimensional reduction}: their effective spectral dimension drops
from 3 toward 2 at large scales when edge-density falls below a critical
threshold. The continuous Laplacian $\nabla^2$ implicitly assumes a 3D-embeddable
graph at all scales. In the sparse halo, this assumption breaks down.

\paragraph{Conditional derivation of Tully-Fisher.}
We define a critical actualization density $\lambda_c$ at which the DAG
transitions from 3D-like (dense core) to 2D-like (sparse halo) behaviour.
The transition radius $r_c$ is the radius at which the Newtonian acceleration
equals the characteristic acceleration $a_0$ associated with $\lambda_c$:
\begin{equation}
  a_N(r_c) = \frac{GM}{r_c^2} = a_0
  \quad\Longrightarrow\quad
  r_c = \sqrt{\frac{GM}{a_0}}.
  \label{eq:critical_radius}
\end{equation}

\textit{If} the causal graph undergoes dimensional reduction at $\lambda_c$
such that outward causal flux is constrained to a boundary scaling as $r$
rather than $r^2$, then causal flux conservation at the transition boundary
requires:
\begin{equation}
  a_T(r) = a_N(r_c)\,\frac{r_c}{r}
           = \frac{GM}{r_c\,r}
           = \frac{\sqrt{GM\,a_0}}{r},
  \label{eq:halo_accel}
\end{equation}
where the second equality substitutes equation~\eqref{eq:critical_radius}.
Setting $v^2/r = a_T(r)$ yields a perfectly flat rotation curve:
\begin{equation}
  v_{\mathrm{flat}}^2 = \sqrt{GM\,a_0}.
  \label{eq:tully_fisher}
\end{equation}
Squaring recovers the baryonic Tully-Fisher relation:
\begin{equation}
  v_{\mathrm{flat}}^4 = a_0\,G\,M.
  \label{eq:tully_fisher_derived}
\end{equation}

This result is now grounded by the Poisson-CSG formula: the RA causal
graph in a region with $\mu(x) = \lambda(x)\,V_{\mathrm{coh}}\,p_{\mathrm{th}}$
has spectral dimension $d_s = 3$ when $\mu \gg 1$ (dense, GR regime)
and $d_s = 2$ when $\mu \ll 1$ (sparse, galactic halo regime), by the
Durhuus-Jonsson-Wheater theorem for directed Lorentzian graphs.
The graph \emph{is} 2-dimensional in galactic halos --- a direct consequence of the
Poisson-CSG topology. Tully-Fisher follows from causal flux conservation
on a locally 2D graph. The empirical MOND acceleration constant
$a_0 \approx 1.2\times10^{-10}$~m/s$^2$ is the macroscopic shadow of
$\lambda_c$ --- the critical graph density where $\mu = 1$, the
Erd\H{o}s-R\'{e}nyi giant-component threshold applied cosmologically.
Hard Wall~2 (Section~\ref{sec:hardwalls}) is now precisely: derive
$p_{\mathrm{th}} = \tau_d/\tau_{\mathrm{mol}}$ from RA-CSG dynamics
to close the identification $\mu = \sigma$ exactly.

% ─────────────────────────────────────────────────────────────────────────────
\section{Gravitational Lensing in Relational Actualism}
\label{sec:lensing}
% ─────────────────────────────────────────────────────────────────────────────

Gravitational lensing is the primary observational constraint on any
dark matter alternative. This section establishes four distinct lensing
regimes in RA, provides a quantitative $\kappa(\theta)$ profile for the
Bullet Cluster, and identifies the open calculations for each regime.

\subsection{Lensing in RA: Two Source Terms and the DJW Regime}
\label{sec:lensing_sources}

In standard GR, the convergence profile is
\begin{equation}
  \kappa(\theta) = \frac{\Sigma_{\mathrm{eff}}(\theta)}{\Sigma_{\mathrm{cr}}},
  \qquad
  \Sigma_{\mathrm{cr}} = \frac{c^2}{4\pi G}\frac{D_S}{D_L D_{LS}}.
\end{equation}
The effective surface density in RA has two contributions,
\begin{equation}
  \Sigma_{\mathrm{eff}}(\mathbf{b})
  = \underbrace{\int \rho_A\,dl}_{\text{accumulated causal depth}}
  + \underbrace{\int \Sigma_\lambda\,dl}_{\text{2D gradient correction}},
  \label{eq:Sigma_RA}
\end{equation}
where $\rho_A$ is the $A_{\mathrm{RA}}$-weighted baryon density and
$\Sigma_\lambda = \xi/(4\pi G R r_d)$ is the disk-geometry correction
(Section~\ref{sec:rotation}).

\paragraph{The DJW lensing regime: an RA-native derivation.}
In the sparse galactic halo ($\mu < 1$), the Durhuus-Jonsson-Wheater theorem
establishes that the Poisson-CSG has effective spectral dimension $d_s = 2$.
In 2D gravity the Poisson equation gives a logarithmic potential; the unique
potential consistent with a flat rotation curve ($v_\mathrm{flat} = $ const) is:
\begin{equation}
  \Phi_{\mathrm{DJW}}(r) = v_{\mathrm{flat}}^2\,\ln r + C.
  \label{eq:log_potential}
\end{equation}
For lensing, null geodesics in the weak-field RA metric satisfy the standard
deflection formula (valid since RA satisfies linearised GR in this regime):
\begin{equation}
  \hat\alpha
  = -\frac{2}{c^2}\int_{-\infty}^{\infty}\frac{\partial\Phi}{\partial b}\,dz.
  \label{eq:deflection}
\end{equation}
Substituting $\Phi_{\mathrm{DJW}} = v_{\mathrm{flat}}^2\ln\sqrt{b^2+z^2}$:
\begin{equation}
  \frac{\partial\Phi}{\partial b} = v_{\mathrm{flat}}^2\,\frac{b}{b^2+z^2},
  \qquad
  \int_{-\infty}^{\infty}\frac{b\,dz}{b^2+z^2} = \pi,
\end{equation}
so that
\begin{equation}
  \hat\alpha = -\frac{2\pi v_{\mathrm{flat}}^2}{c^2}
  \quad\text{[constant, independent of }b\text{].}
  \label{eq:alpha_hat}
\end{equation}
A constant deflection angle is the defining lensing signature of the
$\kappa\propto 1/b$ convergence profile. From equation~\eqref{eq:alpha_hat}:
\begin{equation}
  \boxed{
  \kappa_{\mathrm{DJW}}(b)
  = \frac{v_{\mathrm{flat}}^2}{4G\,\Sigma_{\mathrm{cr}}\,b},
  \qquad
  \theta_E = \frac{2\pi v_{\mathrm{flat}}^2}{c^2}\frac{D_{LS}}{D_S}.
  }
  \label{eq:kappa_DJW}
\end{equation}
This result is \emph{RA-native}: neither the SIS nor any GR assumption was
invoked. The logarithmic potential $\Phi \propto \ln r$ is forced by the DJW
spectral dimension theorem applied to the Poisson-CSG at $\mu < 1$; the
lensing formula follows from the RA weak-field metric. The SIS profile emerges
as a theorem of the DJW causal structure.

\paragraph{Critical distinction from $\Lambda$CDM.}
$\Lambda$CDM realises $\kappa \propto 1/b$ via an NFW dark-matter halo
extending to $\sim$Mpc. RA realises the same shape via DJW-modified baryonic
gravity, cutting off at the baryonic edge ($R_\mathrm{baryon}\sim 150$--300~kpc).
The same $\kappa(b)$ profile shape but different radial extent is directly
testable at $b > 300$~kpc with current survey data.

\subsection{The Bullet Cluster: \texorpdfstring{$\kappa(\theta)$}{kappa(theta)} Profile}
\label{sec:bullet}

The Bullet Cluster (1E~0657-558) is the strongest lensing constraint on
dark matter alternatives~\citep{Clowe2006}. Two galaxy clusters collided;
the baryonic gas shock-heated and decelerated at the collision centre while
the stellar component passed through collisionlessly. Gravitational lensing
is offset $\sim 720$~kpc from the gas peak, co-moving with the stellar
component at $8\sigma$ significance.

\paragraph{The resolution: $A_{\mathrm{RA}}$ vs $\lambda$.}
The RA metric is sourced by $\Pact[T^{\mu\nu}]$ integrated over the full
past light cone $J^-(x)$: accumulated causal depth $A_{\mathrm{RA}}$, not
the instantaneous rate $\lambda(x,t)$. The ratio:
\begin{equation}
  \frac{A_{\mathrm{RA}}(\mathrm{stellar})}{A_{\mathrm{RA}}(\mathrm{gas})}
  \;\approx\;
  \frac{\Gamma_{\mathrm{stellar}}\,T_{\mathrm{stellar}}}
       {\Gamma_{\mathrm{shock}}\,T_{\mathrm{shock}}}
  \;\approx\;
  \frac{10^{14}\,\mathrm{s}^{-1}\times 10\,\mathrm{Gyr}}
       {3\times10^6\,\mathrm{s}^{-1}\times 150\,\mathrm{Myr}}
  \;\approx\; 9.7\times10^7.
  \label{eq:RA_ratio}
\end{equation}
The observed lensing offset requires the stellar component to exceed
the gas by a factor of $\sim 10$; RA provides $\sim 10^8$ headroom.

\paragraph{Why lensing anchors to the galaxies, not the gas.}
This is the point modified-gravity alternatives typically fail to explain.
In RA the answer is structural. The gas-gas collision generates a
high-entropy shock: thermal energy is deposited isotropically, the
BDG action of the shocked gas vertices is large but
\emph{directionally incoherent} --- the causal chains point in many
directions and partially cancel. The resulting contribution to
$A_{\mathrm{RA}}$ per unit mass is low because the BDG path weight
$W = \sum_k c_k N_k$ sums over a thermally randomised distribution
of $N$-vectors with no preferred direction. Entropy is high; directed
causal depth is not.

The collisionless stellar component, by contrast, maintains
\emph{coherent, directional} causal chains accumulated over billions
of years of low-entropy gravitational dynamics. Every star continues
its smooth orbital actualization through the collision, unperturbed.
The stellar $A_{\mathrm{RA}}$ is therefore dominated by long,
coherent directed chains with high BDG path weights. The gas cloud
cannot accumulate comparable $A_{\mathrm{RA}}$ during the
$\sim 150$~Myr shock crossing time, regardless of its instantaneous
energy density.

Formally: the actualization projector $\mathcal{P}_{\mathrm{act}}$
selects ordered, entropy-producing events. Thermal equilibration (the
shocked gas) is entropy-maximising but \emph{not} causal-depth
maximising. Inertial orbital motion (the stars) is causal-depth
maximising. $A_{\mathrm{RA}} \propto \int \mathrm{d}\tau\,\Gamma_{\mathrm{structured}}$
where $\Gamma_{\mathrm{structured}}$ counts only structured, directional
actualization events, not thermal scattering. This is why lensing
traces the stellar component: it is tracking coherent causal depth,
not total energy density.

\paragraph{Quantitative $\kappa(\theta)$ profile.}
Modelling each sub-cluster as a Plummer sphere with mass
$M_{\mathrm{stellar}} = 2\times10^{13}\,M_\odot$ and scale radius
$b_\star = 300$~kpc, and using the $A_{\mathrm{RA}}$-weighted gas
contribution (suppressed by the factor~\eqref{eq:RA_ratio}),
the RA convergence profile is:
\begin{equation}
  \kappa_{\mathrm{RA}}(b)
  = \frac{1}{\Sigma_{\mathrm{cr}}}
    \left[\Sigma_\star(b) + \frac{\Sigma_{\mathrm{gas}}(b)}{A_{\mathrm{ratio}}}\right],
\end{equation}
where $\Sigma_\star$, $\Sigma_{\mathrm{gas}}$ are Plummer surface densities.
The result at representative impact parameters is:

\begin{center}
\renewcommand{\arraystretch}{1.20}
\begin{tabular}{cccccc}
\toprule
$b$ (kpc) & $\Sigma_\star$ ($M_\odot/\mathrm{kpc}^2$)
          & $\kappa_{\mathrm{RA}}$ & Clowe et al.\ $\kappa$ & Ratio \\
\midrule
100 & $5.7\times10^7$ & $0.019$ & $0.35$ & $0.054$ \\
300 & $1.8\times10^7$ & $0.006$ & $0.12$ & $0.050$ \\
500 & $5.0\times10^6$ & $0.002$ & $0.07$ & $0.023$ \\
720 & $1.6\times10^6$ & $0.0005$ & $0.04$ & $0.012$ \\
\bottomrule
\end{tabular}
\end{center}

\noindent
The stellar component overwhelmingly dominates the gas contribution
($\kappa_{\mathrm{gas}} / \kappa_{\mathrm{stellar}} \lesssim 10^{-7}$),
confirming the qualitative result: \emph{lensing traces the stellar
distribution, not the gas}. This is consistent with Clowe et al.~\citep{Clowe2006}.

The amplitude is currently underpredicted by a factor of $\sim 20$ when
only the Plummer stellar contribution is used. This is expected: the $\kappa$
table above uses only instantaneous stellar flux as the $A_{\mathrm{RA}}$ proxy.
The full RA prediction requires two additional contributions that together
substantially reduce the deficit:
\begin{enumerate}[noitemsep]
  \item The $A_{\mathrm{RA}}$-weighted metric source integrates over the
        \emph{full actualization history} of the stellar sub-cluster
        ($\sim 10$~Gyr at $\Gamma_{\mathrm{stellar}} \sim 10^{14}\,\mathrm{s}^{-1}$).
        With realistic stellar population synthesis, the effective $A_{\mathrm{RA}}$
        amplitude increases significantly over the single-rate estimate above.
  \item The hot intracluster medium (ICM, $T \sim 10^8$~K,
        $\Gamma_{\mathrm{ICM}} \sim 10^6\,\mathrm{s}^{-1}$,
        mass $\sim 5\times M_{\mathrm{stellar}}$) contributes a comparable
        $A_{\mathrm{RA}}$ budget to the stars at cluster scales.
\end{enumerate}
The Weyl curvature mapping from $A_{\mathrm{RA}}$ to null geodesic deflection
is required for quantitative verification (Hard Wall~5 below). The
$10^8$ headroom in the $A_{\mathrm{RA}}$ stellar-to-gas ratio
(equation~\eqref{eq:RA_ratio}) is robust to $\mathcal{O}(1)$ corrections.

\paragraph{Gravitational wake.}
Since $A_{\mathrm{RA}}$ accumulates over the stellar lifetime, one must verify
that the stellar trajectory does not produce a Weyl-curvature wake inconsistent
with the observed profile. RA does not predict such a wake: $A_{\mathrm{RA}}$
encodes causal depth at the \emph{current} graph state, not a convolution of
historical positions. The DAG records what has actualized, not where it has been.
Formalising this in the covariant continuum limit is Hard Wall~5.

\subsection{Lensing Regimes: A Classification}
\label{sec:lensing_regimes}

Gravitational lensing in RA operates differently across four distinct
physical regimes, which we classify here for the first time.

\paragraph{Regime 1: Isolated galaxies.}
For isolated galaxies the primary lensing mechanism is the
\emph{DJW dimensional-reduction regime} (Section~\ref{sec:dimensional_reduction}).
At radii $r > r_c$ where the actualization density $\mu < 1$, the
causal graph is locally 2-dimensional by the Durhuus-Jonsson-Wheater theorem,
and the gravitational force law becomes $F \propto 1/r$ rather than $1/r^2$.
This has the same $\kappa \propto 1/b$ convergence profile as a Singular Isothermal Sphere (SIS), but with half the deflection angle: the RA-native derivation (equation~\eqref{eq:kappa_DJW}) gives $\hat\alpha_{\mathrm{RA}} = 2\pi v^2/c^2$ vs.\ $\hat\alpha_{\mathrm{SIS}} = 4\pi\sigma^2/c^2$
within the baryon distribution, giving:
\begin{equation}
  \kappa_{\mathrm{RA}}(b) = \frac{v_{\mathrm{flat}}^2}{4 G\,\Sigma_{\mathrm{cr}}\,b},
  \qquad
  \theta_E = \frac{2\pi v_{\mathrm{flat}}^2}{c^2}\frac{D_{LS}}{D_S}.
  \label{eq:SIS_kappa}
\end{equation}
For a Milky-Way analog ($v_{\mathrm{flat}} = 220$~km~s$^{-1}$,
$D_L = 500$~Mpc, $D_S = 1500$~Mpc), this gives an Einstein radius
$\theta_E = 0.47$~arcsec for a Milky-Way-analog with $v_\mathrm{flat} = 220$~km~s$^{-1}$
(rising to 0.8--1.2~arcsec for more massive early-type lenses with
$v_\mathrm{flat} \approx 310$--380~km~s$^{-1}$~\citep{Koopmans2009}),
and convergence profile:
\begin{center}
\renewcommand{\arraystretch}{1.15}
\begin{tabular}{cccc}
\toprule
$b$ (kpc) & $\kappa_{\mathrm{DJW}}$ (eq.\,\ref{eq:kappa_DJW}) & $\kappa_{\mathrm{obs}}$ & Ratio \\
\midrule
10  & 0.036 & $0.12$ & 0.30 \\
50  & 0.007 & $0.04$ & 0.18 \\
100 & 0.006 & $0.025$ & 0.22 \\
200 & 0.003 & $0.013$ & 0.21 \\
500 & 0.001 & $0.006$ & 0.19 \\
\bottomrule
\end{tabular}
\end{center}
At $b = 10$~kpc the RA DJW prediction gives $\kappa = 0.036$
(observed $\sim 0.12$, a factor $\sim 3$). At $b = 100$~kpc: $\kappa = 0.006$
(observed $\sim 0.025$, factor $\sim 4$). The deficit is consistent at all radii.
This residual is accounted for by the warm circumgalactic medium
(CGM, $T \sim 10^5$--$10^6$~K, mass $\sim 10^{10}\,M_\odot$, extent
$\sim 150$~kpc~\citep{Tumlinson2017}): the CGM actualizes thermally
($\Gamma_{\mathrm{CGM}} \sim n_e\sigma_T c$), extends the DJW-active
baryon distribution to larger radii, and contributes additional lensing.

\textbf{Key falsifiable prediction:} RA predicts the galaxy-galaxy
lensing signal drops off exponentially at $b > R_{\mathrm{CGM}} \sim 150$--300~kpc,
following the baryon edge. $\Lambda$CDM's NFW halos extend to $\sim$Mpc scales.
The profile shape at $b > 300$~kpc directly discriminates the two frameworks.
\textbf{Status: qualitative prediction complete; quantitative computation
of the full baryon+DJW $\kappa(b)$ profile is the remaining open calculation.}

\paragraph{Regime 2: Colliding galaxies / Bullet Cluster.}
The $A_{\mathrm{RA}}$ ratio (stellar/gas) of $\sim 10^8$ explains
why the lensing offset follows the stellar component. This is RA's
cleanest lensing result, independent of the warm baryon question.
The amplitude calibration requires the ICM contribution and
Weyl mapping (Hard Wall~5). \textbf{Status: qualitatively complete,
quantitatively open.}

\paragraph{Regime 3: Galaxy clusters.}
The hot ICM ($T \sim 10^8$~K) has $\Gamma_{\mathrm{ICM}} \approx 10^6$~s$^{-1}$
per baryon and mass $\sim 5\times M_{\mathrm{stellar}}$. Its total $A_{\mathrm{RA}}$
contribution per unit mass is $\sim 10^{-8}$ that of stars, but the ICM budget
may bring total RA lensing to within a factor $\sim 2$--$3$ of observations
without dark matter. Additionally, proximity and tidal effects during cluster
assembly enhance $A_{\mathrm{RA}}$ in the outer halo relative to isolated
galaxies. \textbf{Status: order-of-magnitude estimate, quantitative
computation needed.}

\paragraph{Regime 4: Large-scale structure and cosmic shear.}
At scales $\gg$ the void correlation length ($\sim 100$~Mpc), RA approaches
the GR limit ($\Pact \approx \mathrm{id}$) and the lensing power spectrum
recovers $\Lambda$CDM to within $\sim 1\%$. At intermediate scales (10--100~Mpc),
RA corrections from the actualization density field become detectable in
principle. The lensing power spectrum in RA requires running CLASS or CAMB
with the RA source term $G_{\mu\nu} = 8\pi G\,\Pact[T_{\mu\nu}]$.
\textbf{Status: open calculation, awaits the cosmological simulation (D2).}

\subsection{Hard Wall: Weyl Curvature Sourcing}
\label{sec:HW_weyl}

The common technical gap across all four regimes is the mapping from the
RA metric source $G_{\mu\nu} = 8\pi G\,\Pact[T_{\mu\nu}]$ to the
Weyl tensor $C_{\mu\nu\rho\sigma}$ that governs null geodesic deflection.

In vacuum regions between mass concentrations, $G_{\mu\nu} = 0$
(no actualized matter), but $C_{\mu\nu\rho\sigma} \neq 0$:
the Weyl curvature there is determined by the boundary conditions
from the surrounding $A_{\mathrm{RA}}$-weighted mass distribution,
propagated via the vacuum Einstein equations. This is a standard GR
boundary-value problem once the RA source at mass locations is established.

The $10^8$ headroom in equation~\eqref{eq:RA_ratio} is robust to
$\mathcal{O}(1)$ geometric factors in the Weyl propagation.
The precise $\kappa(\theta)$ map requires:
\begin{enumerate}[noitemsep]
  \item The full covariant RA field equations from RACL~\citep{Sandeman2026RACL};
  \item A realistic $A_{\mathrm{RA}}(\mathbf{x})$ field from stellar
        population synthesis;
  \item A numerical weak-field propagation of the RA metric perturbation
        to produce $\kappa(\theta)$ maps comparable to Clowe et al.
\end{enumerate}
This is the primary remaining hard wall for RADM.
\textbf{Collaborator target:} gravitational lensing observers;
weak-field numerical GR; stellar population synthesis.

% ─────────────────────────────────────────────────────────────────────────────


The Bullet Cluster (1E~0657-558) is the strongest observational constraint on dark
matter alternatives~\citep{Clowe2006}. Two galaxy clusters collided. The baryonic
gas shock-heated ($T \sim 10^8$~K) and decelerated at the collision centre. The
stellar component passed through collisionlessly. Gravitational lensing is offset
$\sim 720$~kpc from the gas, co-moving with the stellar component (8$\sigma$
significance).

\paragraph{The naive failure.}
If the metric tracked instantaneous $\lambda$, the shocked gas---with
$\Gamma_{\mathrm{shock}} \sim 3\times\Gamma_{\mathrm{pre}}$---would dominate
lensing. This contradicts observation.

\paragraph{The resolution: $A_{\mathrm{RA}}$ vs $\lambda$.}
The RAGC metric is sourced by $\Pact[T^{\mu\nu}]$ integrated over the full past
light cone $J^-(x)$---it is the accumulated causal depth $\ARA$, not the
instantaneous rate $\lambda(x,t)$, that sources the background metric. The
stellar component has actualized at $\Gamma_{\mathrm{stellar}} \sim
10^{14}$~s$^{-1}$ per baryon for $\sim 10$~Gyr. The gas has actualized at
$\Gamma_{\mathrm{gas}} \sim 10^6$~s$^{-1}$ per baryon. The accumulated
causal depth ratio:
\begin{equation}
  \frac{\ARA(\mathrm{stellar})}{\ARA(\mathrm{gas})}
  \;\approx\; \frac{\Gamma_{\mathrm{stellar}}}{\Gamma_{\mathrm{gas}}}
  \;\approx\; 10^8.
  \label{eq:RA_ratio_simplified}
\end{equation}
The Bullet Cluster lensing requires the stellar component to exceed the gas
contribution by a factor of $\sim 10$. The $\ARA$ ratio provides $10^8$
orders of magnitude of headroom. The 150~Myr shock contributes $\sim 4.5$\% of
the gas's total accumulated causal depth---a negligible perturbation to the
background metric. No new mechanism. No new parameters. The same distinction
that makes biological assembly depth meaningful in RACI~\citep{Sandeman2026c}
resolves the Bullet Cluster in RADM.

\paragraph{Weyl curvature and the lensing signal.}
A precise technical concern requires acknowledgement: gravitational lensing traces
null geodesics, which are governed by the Weyl tensor $C_{\mu\nu\rho\sigma}$
(the trace-free part of the Riemann tensor), whereas the RA-modified field
equations source the Ricci tensor $R_{\mu\nu}$ via the local stress-energy.
The Bullet Cluster argument above establishes that $\ARA$ dominates the metric
source over astrophysically relevant timescales. However, demonstrating explicitly
how the accumulated $\ARA$ field contributes to the Weyl curvature at the moment
of collision --- and confirming that the 10~Gyr stellar trajectory does not
produce a gravitational wake inconsistent with the observed lensing profile ---
requires the full covariant field equations (RAGC Derivation~3,
Hard Wall~iv~\citep{Sandeman2026b}).

In the weak-field limit, the Weyl and Ricci contributions are not independent:
for a quasi-static mass distribution, Weyl curvature is sourced by the tidal
field of the stress-energy distribution, and a dominant $\ARA$ contribution to
$T^{\mu\nu}$ will produce a corresponding dominant contribution to the Weyl
curvature. The $10^8$ headroom in equation~\eqref{eq:RA_ratio} suggests this
remains qualitatively robust. The explicit covariant mapping is Hard Wall~5
below.

\paragraph{The gravitational wake concern.}
A referee concern deserving explicit acknowledgement: because $A_{\mathrm{RA}}$
accumulates over the stellar lifetime ($\sim 10$~Gyr), one must verify that
the integrated stellar trajectory does not produce a gravitational wake --- a
trailing Weyl-curvature signature inconsistent with the observed lensing
profile. The RA framework does not predict such a wake for the following
reason: $A_{\mathrm{RA}}$ encodes causal depth at the \emph{current} graph
state of the stellar component, not a convolution of historical positions.
The causal DAG records \emph{what has actualized}, not \emph{where it has
been}. Stellar trajectories write vertices at their instantaneous boundary;
the graph does not retain a spatially extended record of the trajectory.
Formalising this argument --- showing that the DAG metric update does not
produce a wake in the continuum limit --- is part of Hard Wall~5.

\paragraph{Hard wall: quantitative verification.}
The order-of-magnitude argument is encouraging but not a proof. Computing the
stellar-to-gas $\ARA$ ratio with realistic stellar population synthesis models,
and mapping $\ARA$ to lensing convergence $\kappa(\theta)$ via the full covariant
Weyl sourcing, is Hard Wall~3 (Section~\ref{sec:hardwalls}).

% ─────────────────────────────────────────────────────────────────────────────
\section{Open Problems and Hard Walls}

\textbf{Update (RACL):} The continuum-limit hard wall shared with
RAGC~\citep{Sandeman2026b}---proving that the Lean-verified Local
Ledger Condition uniquely implies the Einstein field equations---is now
closed by RACL~\citep{Sandeman2026RACL}. The field equation
$G_{\mu\nu} = 8\pi G\,\mathcal{P}_{\mathrm{act}}[T_{\mu\nu}]$ with
$\Lambda = 0$ is proved as a theorem with no free parameters.
The field equation RADM uses as its starting point is therefore derived,
not assumed. The remaining hard walls below are specific to the
RADM programme.

\label{sec:hardwalls}
% ─────────────────────────────────────────────────────────────────────────────

\begin{description}
  \item[Hard Wall 1: Microscopic Derivation of $\xi$ and the Jacobson Coefficient.]
  The coupling constant is numerically observed to be
  $\xi/(G M_{\mathrm{disk}}) \approx 0.28$ for a Milky Way analog.
  RAGC~\citep{Sandeman2026b} establishes via a thermodynamic derivation that
  $\xi$ must emerge from the ratio of microscopic actualization entropy to the
  fundamental RA area element:
  \begin{equation}
    \frac{\xi}{G M_{\mathrm{disk}}}
      = \alpha \times 4\!\left(\frac{l_P}{l_{\mathrm{RA}}}\right)^{\!2}
        \langle\Delta S_{\mathrm{vertex}}\rangle,
    \label{eq:xi_consistency}
  \end{equation}
  where $\alpha$ is an $\mathcal{O}(1)$ coefficient to be derived from the
  discrete Jacobson calculation, and $\langle\Delta S_{\mathrm{vertex}}\rangle$
  is the mean relative entropy increase per actualization event.

  A quantitative consistency check with the Kinematic Coherence Bound of
  RAQI~\citep{Sandeman2026d} yields a striking result. Setting
  $l_{\mathrm{RA}} = l_P$ and $\langle\Delta S_{\mathrm{vertex}}\rangle
  = p_{\mathrm{th}} \approx 10^{-2}$ (the fault-tolerance threshold):
  \begin{equation}
    \alpha_{\mathrm{required}} = \frac{\xi/GM_{\mathrm{disk}}}{4\,p_{\mathrm{th}}}.
    \label{eq:alpha_required}
  \end{equation}
  The Jacobson derivation applied to the discrete actualization sector gives
  $\alpha = 2\pi \approx 6.28$ (the Unruh temperature periodicity factor).
  The value of $\alpha_{\mathrm{required}}$ depends critically on which
  components of the baryonic disk are included in $M_{\mathrm{disk}}$.

  Using the stellar-disk-only mass $M = 6.0\times10^{10}\,M_\odot$ yields
  $\alpha_{\mathrm{required}} \approx 7.0$ --- a 12\% gap above $2\pi$.
  However, this comparison is known to underestimate the true disk mass.
  Including the gaseous disk (${\approx}0.8\times10^{10}\,M_\odot$) and
  thick disk (${\approx}1.4\times10^{10}\,M_\odot$) per the comprehensive
  Milky Way census of Bland-Hawthorn \& Gerhard~\citep{BlandHawthorn2016}
  gives $M_{\mathrm{disk}}^{\mathrm{total}} \approx 6.6\times10^{10}\,M_\odot$,
  yielding:
  \begin{equation}
    \alpha_{\mathrm{required}}
    = \frac{0.256}{4\times 10^{-2}} \approx 6.39
    \approx 2\pi + 1.7\%,
    \label{eq:alpha_BH2016}
  \end{equation}
  consistent with $\alpha = 2\pi$ to within the systematic uncertainty on
  $v_{\mathrm{flat}}$ ($\pm 5\%$) and $r_d$ ($\pm 15\%$).
  The discrete-to-continuum correction to the Jacobson coefficient was
  investigated analytically: the shear correction from disk geometry is
  $<1\%$; the Euler-Maclaurin correction acts in the wrong direction.
  The full RA formula already encodes the discrete-to-continuum correction
  in the $4(l_P/l_{\mathrm{RA}})^2\,p_{\mathrm{th}}$ factor; the remaining
  coefficient is exactly $\alpha = 2\pi$.
  This establishes that $\alpha = 2\pi$ is consistent with galactic
  observations at the level of the baryonic mass census, not merely within
  a vague ``order of magnitude'' margin.
  \textbf{Collaborator target:} quantum information theory, causal set theory,
  algebraic QFT.

  \item[Hard Wall 2: Formal Proof of Dimensional Reduction at $\lambda_c$.]
  Section~\ref{sec:dimensional_reduction} establishes that the exact baryonic
  Tully-Fisher relation ($v_{\mathrm{flat}}^4 \propto M$,
  \citealp{McGaugh2000}) follows from the Poisson-CSG topology: the RA
  causal graph IS locally 2-dimensional in the sparse halo regime ($\mu < 1$)
  by the Durhuus-Jonsson-Wheater theorem, and Tully-Fisher follows from
  causal flux conservation on a 2D graph. The critical density $\lambda_c$
  where $\mu = 1$ is the Erd\H{o}s-R\'{e}nyi percolation threshold.
  Hard Wall~2 is now precisely: derive $p_{\mathrm{th}} = \tau_d/\tau_{\mathrm{mol}}$
  from RA-CSG dynamics, and run the BDG-filter MCMC to extract
  the exact $\lambda_c$ prediction. The empirical MOND acceleration constant
  $a_0 \approx 1.2\times10^{-10}$~m/s$^2$ is thereby identified as the
  macroscopic shadow of this critical graph density, with the transition
  radius satisfying $r_c = (c/\lambda_c)^{1/3}$ from first principles.

  While generic critical trees exhibit anomalous diffusion with spectral
  dimension $d_s = 4/3$ (Alexander-Orbach conjecture), the
  Durhuus-Jonsson-Wheater theorem rigorously establishes that for
  Lorentzian causal triangulations, the directed time-like structure biases
  the random walk and yields exactly $d_s = 2$ in the sparse
  branched-polymer phase. The directed edges of the RA causal DAG forbid the
  random walker from diffusing backward into dead-end spatial branches,
  confining it to a quasi-1D time-like trajectory with spatial leaves. It is
  this causal structure that forces $d_s = 2$ rather than $4/3$.

  The precise remaining gap is to prove that the RA-CSG stochastic
  graph-rewriting rules of the Engine of
  Becoming~\citep{Sandeman2026RAEB} fall into the same Lorentzian causal
  triangulations universality class. Concretely: perform Markov Chain Monte
  Carlo simulations of the RA actualization dynamics to calculate $d_s$ as a
  function of $\lambda$, verify the transition from $d_s = 3$ (dense core)
  to $d_s = 2$ (sparse halo) at the threshold $\lambda_c \sim c/r_c^3$, and
  determine the exact critical exponent $\beta$ for the transition sharpness.
  The galaxy-to-galaxy variation in $\xi/(GM_{\mathrm{disk}})$ (Section~3.1)
  is explained once $\lambda_c$ is computed locally for each galaxy's disk
  boundary. The Poisson-CSG formula of RAGC~\citep{Sandeman2026b} provides the
  mechanism: $c_k^{\mathrm{RA}}(x) = e^{-\mu}\mu^k/k!$ with
  $\mu = \lambda V_{\mathrm{coh}} p_{\mathrm{th}}$ makes the spectral
  dimension local --- $d_s \to 3$ as $\mu \gg 1$ and $d_s \to 2$ as
  $\mu \ll 1$ --- with the transition at $\mu = 1$ defining $\lambda_c$.
  The exact $c_k$ beyond the mean-field limit requires the
  BDG-filter MCMC of RAGC Derivation~1. Testable against
  SPARC~\citep{Lelli2016}.
  \textbf{Collaborator target:} causal set theory, CDT computational
  physics, spectral graph theory.

  \item[Hard Wall 3: Quantitative Bullet Cluster.] Computing the stellar-to-gas
  $\ARA$ ratio with realistic stellar population synthesis models; mapping $\ARA$
  to lensing convergence $\kappa(\theta)$ and marginalizing over stellar
  mass-to-light ratios and gas profiles. Decisive test: reproducing the $8\sigma$
  lensing offset amplitude and spatial profile.
  \textbf{Collaborator target:} observational astrophysics, stellar evolution.

  \item[Hard Wall 4: CMB acoustic structure and early structure formation.]
  In $\Lambda$CDM, collisionless CDM seeds early structure formation and
  governs CMB acoustic peak ratios.  For RADM to constitute a complete dark
  matter replacement, the topological boundary mechanism must produce
  sufficient early density perturbation amplification, and the RA-modified
  field equations must reproduce the observed CMB acoustic peak structure
  without a CDM parameter.

  \textbf{New partial result.}
  For a baryon perturbation $\delta_b \propto \cos(k c_s t)$, the
  actualization-density perturbation satisfies:
  \begin{equation}
    \frac{\delta_A}{\delta_b} = \frac{\sin(k r_s)}{k r_s},
    \label{eq:sinc}
  \end{equation}
  where $r_s$ is the sound horizon.  At every CMB acoustic peak
  ($k r_s = n\pi$), $\sin(n\pi)/(n\pi) = 0$ exactly: the $\rho_A$
  perturbation vanishes at the peaks regardless of the relaxation timescale
  $\tau_{\mathrm{relax}}$, $f_0$, or any other parameter.  The actualization
  density is therefore CDM-like \emph{in shape} at all acoustic peaks --- a
  structural match requiring no tuning.

  \textbf{New cosmological derivation (RASM/ra\_bdg\_couplings).}
  The dark matter/baryon density ratio is derived from BDG combinatorics:
  \begin{equation}
    \frac{\Omega_{\mathrm{CDM}}}{\Omega_b}
    = \frac{W_{\mathrm{other}}}{W_{\mathrm{baryon}}} \times \alpha_s(2m_p)
    = 17.32 \times 0.312 = 5.40,
    \label{eq:f0_from_RASM}
  \end{equation}
  where $W_{\mathrm{baryon}} = 3/2$ exactly and $W_{\mathrm{other}}/W_{\mathrm{baryon}} = 17.32$
  follow from the BDG integers $(1,-1,9,-16,8)$ at $\mu=1$, and
  $\alpha_s(2m_p) = 0.312 \pm 0.004$ is the strong coupling at the nuclear
  binding scale (FLAG 2021 lattice + 2-loop QCD running).  This agrees with
  Planck~2018 $\Omega_{\mathrm{CDM}}/\Omega_b = 5.416$ to $0.3\%$
  \citep{Sandeman2026RASM,Sandeman2026BDG}.

  The full CMB peak amplitude problem --- showing that the $\rho_A$ amplitude
  equals the observed $\sim 5\times$ baryon enhancement --- follows from
  equation~\eqref{eq:f0_from_RASM} and is now structurally closed.
  The remaining targets are: (i)~derive the RA renormalization group
  $d\alpha_s/d\ln\rho$ from first principles (removing the QCD running
  as external input); and (ii)~compute the transfer function and CMB
  power spectrum from the RA perturbation equations.
  \textbf{Collaborator target:} early-universe cosmology; CMB phenomenology;
  cosmological perturbation theory.

  \item[Hard Wall 5: Weyl curvature sourcing and the Bullet Cluster lensing
  prediction.] Gravitational lensing traces the Weyl tensor
  $C_{\mu\nu\rho\sigma}$ via null geodesics, while the RA-modified field
  equations source the Ricci tensor $R_{\mu\nu}$ via local actualization
  density. Three RA-internal targets remain:
  \begin{enumerate}[noitemsep]
  \item \textbf{Covariant field equation [resolved in RAGC].}
  The covariant form $\Box(\ln\lambda) = g^{\mu\nu}\nabla_\mu\nabla_\nu(\ln\lambda)$
  is proved (RAGC Proposition~9, Corollary~10~\citep{Sandeman2026b}):
  $\Box(\ln\lambda) = \nabla^2(\ln\lambda) + O(\Phi)$ in the static
  weak-field limit, and the modified Poisson equation
  $\nabla^2\Phi = 4\pi G\rho_A + \xi\,\nabla^2(\ln\lambda)$ is derived.
  The $\lambda$ equation of motion $\nabla^2(\ln\lambda) = -\rho_A/(3\xi)$
  is established, predicting the dark-matter halo profile from first principles.
  \item \textbf{Weyl sourcing from $A_{\mathrm{RA}}$ [open].}
  Demonstrating quantitatively how accumulated causal depth $\ARA$
  generates the observed $8\sigma$ lensing offset between the
  gas centroid and the lensing peak in the Bullet Cluster.
  This requires mapping $\ARA$ to the Weyl potential; the tidal
  decomposition gives $\Psi - \Phi \propto \nabla^{-2}(\rho_\lambda - \bar\rho_\lambda)$,
  and the $10^8$ stellar/gas $\ARA$ headroom (equation~\eqref{eq:RA_ratio})
  is expected to survive. Precise quantitative prediction outstanding.
  \item \textbf{No-wake prediction [open].}
  The DAG records current causal depth, not historical spatial trajectories;
  therefore no gravitational wake from past stellar positions should appear
  in the Weyl curvature after a cluster collision. This is a
  \emph{novel falsifiable prediction of RA} distinguishing it from
  any particle dark matter model. Quantifying the wake amplitude
  requires the covariant $\lambda$ profile at cluster-collision timescales.
  \end{enumerate}
  \textbf{Primary open target:} derive $\xi$ from RA microscopic structure
  (RAGC Derivation~6~\citep{Sandeman2026b}), then compute the Weyl potential
  from the predicted $\lambda$ profile.
  \textbf{Collaborator targets:} gravitational lensing, cluster astrophysics,
  large-scale structure.
\end{description}

% ─────────────────────────────────────────────────────────────────────────────
\section{Relationship to the RA Suite}
\label{sec:suite}
% ─────────────────────────────────────────────────────────────────────────────

RADM is the sixth companion paper in the RA suite. The logical dependencies are:

\begin{description}
  \item[\textbf{RAQM}~\citep{Sandeman2026a}] Provides the foundational
    actualization criterion and the Kinematic Snap.
  \item[\textbf{RAGC}~\citep{Sandeman2026b}] Establishes the WIMP Prohibition
    (Prediction~4) and the WEP recovery argument (Derivation~6). RADM depends
    on these results.
  \item[\textbf{RACI}~\citep{Sandeman2026c}] The $\ARA$/\,$\lambda$ distinction
    used in the Bullet Cluster resolution is established in RACI. The Bullet
    Cluster is the gravitational analogue of the biological assembly depth result.
  \item[\textbf{RAQI}~\citep{Sandeman2026d}] Provides the quantum thermodynamic
    grounding for the Landauer cost derivations that underpin the actualization
    flux calculations.
  \item[\textbf{RAHC}~\citep{Sandeman2026RAHC}] The recursive coarse-graining
    hierarchy and the topological motif direction (Epilogue) are the longer-range
    context for the RADM conjectures.
\end{description}

\noindent
The suite also includes the Standard Model gauge structure paper
(RASM~\citep{Sandeman2026RASM}), the particle classification
paper (RATM~\citep{Sandeman2026RATM}), and the programme overview
(Foundation~\citep{Sandeman2026Foundation}), and the astrobiology
habitability paper (RACF~\citep{Sandeman2026RACF}).


% ─────────────────────────────────────────────────────────────────────────────
\section{Conclusion}
\label{sec:conclusion}
% ─────────────────────────────────────────────────────────────────────────────

RADM proposes that the dark matter signatures are actualization-density effects
rather than evidence for a new particle species. The WIMP Prohibition (established
in RAGC) rules out causally silent candidates. The $\nabla^2(\ln\lambda)$
mechanism produces flat rotation curves at natural coupling for exponential disk
profiles. The $\ARA/\lambda$ distinction resolves the Bullet Cluster lensing
offset without new parameters. Four hard walls are named precisely.

This paper does not claim to have replaced dark matter. It claims to have
identified a research direction, grounded in the RA framework's single ontological
commitment, that may do so once the hard walls are resolved. The Tully-Fisher
scaling and CMB structure formation requirements are the decisive tests. Until
they are met, this paper is a well-motivated conjecture with natural-order coupling
and two qualitative successes---not a completed model.

% ─────────────────────────────────────────────────────────────────────────────

\section*{Declarations}
\textbf{Funding:} Not applicable.

\textbf{Ethics statement:} This article does not contain any studies with human participants or animals performed by the author.

\quad
\textbf{Competing interests:} The author declares no competing interests.

\textbf{Preprint availability:} This manuscript is available as a preprint at
\href{https://doi.org/10.5281/zenodo.19198513}{doi.org/10.5281/zenodo.19198513}.

\textbf{AI use disclosure:} Claude (Anthropic, model \texttt{claude-sonnet-4-6}) was used in the preparation of this manuscript, including drafting and revision of text, mathematical derivation assistance, \LaTeX{} compilation, and Lean~4 proof development. The author is solely responsible for all scientific content, mathematical claims, and conclusions. Gemini (Google DeepMind, Gemini~2.0~Pro) was used for supplementary literature checking and cross-verification of symbolic calculations.


\textbf{Code availability:} The rotation curve Python script and all companion paper source files are publicly available at \href{https://github.com/jsandeman/Relational-Actualism}{github.com/jsandeman/Relational-Actualism}.

\textbf{Data availability:} This paper presents theoretical work. No datasets
were generated or analysed.

\bibliographystyle{plainnat}
\begin{thebibliography}{99}

\bibitem[Koopmans et al.(2009)]{Koopmans2009}
L.~V.~E. Koopmans, T.~Treu, A.~S. Bolton, S.~Burles, and L.~A. Moustakas.
\newblock The Sloan Lens ACS survey. III. The structure and formation of
  early-type galaxies and their evolution since $z \sim 1$.
\newblock \textit{Astrophys.\ J.}, 703:L51--L54, 2009.
\newblock \href{https://doi.org/10.1088/0004-637X/703/1/L51}{doi:10.1088/0004-637X/703/1/L51}.


\bibitem[Sandeman(2026RACL)]{Sandeman2026RACL}
J.~F. Sandeman.
\newblock {Relational Actualism: The Einstein Field Equations as the Unique
  Macroscopic Limit of the Actualization Graph (RACL)}.
\newblock \textit{Preprint}, 2026.
\newblock \href{https://doi.org/10.5281/zenodo.19270020}{doi.org/10.5281/zenodo.19270020}.


\bibitem[Tumlinson et al.(2017)]{Tumlinson2017}
J.~Tumlinson, M.~S. Peeples, and J.~K. Werk.
\newblock The circumgalactic medium.
\newblock \textit{Annu.\ Rev.\ Astron.\ Astrophys.}, 55:389--432, 2017.
\newblock \href{https://doi.org/10.1146/annurev-astro-091916-055240}{doi:10.1146/annurev-astro-091916-055240}.


\bibitem[Sandeman(2026RATM)]{Sandeman2026RATM}
J.~F. Sandeman.
\newblock {Relational Actualism: The Topology of Matter (RATM)}.
\newblock \textit{Preprint}, 2026.
\newblock \href{https://doi.org/10.5281/zenodo.19265651}{doi.org/10.5281/zenodo.19265651}.


\bibitem[Sandeman(2026a)]{Sandeman2026a}
J.~F.~Sandeman.
\newblock Relational Actualism and Quantum Mechanics ({RAQM}).
\newblock \textit{Foundations of Physics} (Under review), 2026.
\newblock \href{https://doi.org/10.5281/zenodo.19174380}{doi.org/10.5281/zenodo.19174380}.

\bibitem[Sandeman(2026b)]{Sandeman2026b}
J.~F.~Sandeman.
\newblock Relational Actualism: Gravity and Cosmology ({RAGC}).
\newblock Preprint, 2026.
\newblock \href{https://doi.org/10.5281/zenodo.19197999}{doi.org/10.5281/zenodo.19197999}.

\bibitem[Sandeman(2026c)]{Sandeman2026c}
J.~F.~Sandeman.
\newblock Relational Actualism: Complexity and Intelligence ({RACI}).
\newblock Preprint, 2026.
\newblock \href{https://doi.org/10.5281/zenodo.19198224}{doi.org/10.5281/zenodo.19198224}.

\bibitem[Sandeman(2026d)]{Sandeman2026d}
J.~F.~Sandeman.
\newblock Relational Actualism: Implications for Quantum Information ({RAQI}).
\newblock Preprint, 2026.
\newblock \href{https://doi.org/10.5281/zenodo.19198285}{doi.org/10.5281/zenodo.19198285}.

\bibitem[Sandeman(2026e)]{Sandeman2026RAHC}
J.~F.~Sandeman.
\newblock Algorithmic Causal Dynamics: A Unified Framework for Emergence and
  Complexity in Relational Actualism ({RAHC}).
\newblock Preprint, 2026.
\newblock \href{https://doi.org/10.5281/zenodo.19198171}{doi.org/10.5281/zenodo.19198171}.

\bibitem[Bland-Hawthorn \& Gerhard(2016)]{BlandHawthorn2016}
J.~Bland-Hawthorn and O.~Gerhard.
\newblock The Galaxy in context: Structural, kinematic, and integrated
  properties.
\newblock \emph{Annual Review of Astronomy and Astrophysics}, 54:529--596, 2016.
\newblock \href{https://doi.org/10.1146/annurev-astro-081915-023441}{doi:10.1146/annurev-astro-081915-023441}.

\bibitem[Sandeman(2026g)]{Sandeman2026RAEB}
J.~F.~Sandeman.
\newblock The Engine of Becoming: A Covariant Stochastic Graph Rewriting
  Algorithm as the Unified Foundation of Quantum Mechanics and General
  Relativity ({EB}).
\newblock Preprint, 2026.
\newblock \href{https://doi.org/10.5281/zenodo.19198120}{doi.org/10.5281/zenodo.19198120}.

\bibitem[Clowe et~al.(2006)]{Clowe2006}
D.~Clowe, M.~Brada\v{c}, A.~H.~Gonzalez, M.~Markevitch, S.~W.~Randall,
C.~Jones, and D.~Zaritsky.
\newblock A direct empirical proof of the existence of dark matter.
\newblock \textit{The Astrophysical Journal Letters}, 648:L109--L113, 2006.

\bibitem[Lelli et~al.(2016)]{Lelli2016}
F.~Lelli, S.~S.~McGaugh, and J.~M.~Schombert.
\newblock {SPARC}: Mass models for 175 disk galaxies.
\newblock \textit{The Astronomical Journal}, 152:157, 2016.

\bibitem[McGaugh et~al.(2000)]{McGaugh2000}
S.~S.~McGaugh, J.~M.~Schombert, G.~D.~Bothun, and W.~J.~G.~de~Blok.
\newblock The baryonic Tully-Fisher relation.
\newblock \textit{The Astrophysical Journal Letters}, 533:L99--L102, 2000.

\bibitem[Milgrom(1983)]{Milgrom1983}
M.~Milgrom.
\newblock A modification of the Newtonian dynamics.
\newblock \textit{The Astrophysical Journal}, 270:365--370, 1983.

\bibitem[Verlinde(2011)]{Verlinde2011}
E.~Verlinde.
\newblock On the origin of gravity and the laws of Newton.
\newblock \textit{Journal of High Energy Physics}, 2011(4):29, 2011.

\bibitem[Sandeman(2026h)]{Sandeman2026RASM}
J.~F.~Sandeman.
\newblock Relational Actualism and the Standard Model ({RASM}).
\newblock Preprint, 2026.
\newblock \href{https://doi.org/10.5281/zenodo.19229335}{doi.org/10.5281/zenodo.19229335}.

\bibitem[Sandeman(2026BDG)]{Sandeman2026BDG}
J.~F.~Sandeman.
\newblock Two Standard Model Coupling Constants from the Unique 4-Dimensional
  BDG Integers.
\newblock Preprint, 2026.
\newblock \href{https://doi.org/10.5281/zenodo.19324790}{doi.org/10.5281/zenodo.19324790}.

\bibitem[Sandeman(2026i)]{Sandeman2026Foundation}
J.~F.~Sandeman.
\newblock Relational Actualism: A Foundation Paper.
\newblock Preprint, 2026.
\newblock \href{https://doi.org/10.5281/zenodo.19229438}{doi.org/10.5281/zenodo.19229438}.


\bibitem[Sandeman(2026RACF)]{Sandeman2026RACF}
J.~F.~Sandeman.
\newblock The Causal Firewall: Substrate-Independent Physical Constraints
  on the Distribution of Life in the Universe (RACF).
\newblock \textit{Int.\ J.\ Astrobiol.} (submitted), 2026.
\newblock \href{https://doi.org/10.5281/zenodo.19319374}%
  {doi:10.5281/zenodo.19319374}.

\end{thebibliography}

\end{document}
